\section{MLC LLM: la implementación en el dispositivo en la era de los modelos grandes}

\subsection{De TVM a MLC LLM}

MLC LLM es una aplicación vertical de TVM------un motor de inferencia de LLM multiplataforma capaz de compilar e implementar modelos de lenguaje grandes en diversos dispositivos como iPhone, Android, tarjetas gráficas de AMD y MacBook. Utiliza tecnología de compilación de aprendizaje automático, porque el equipo «sabe con toda claridad que no tenemos la capacidad suficiente para extender manualmente el soporte a cada plataforma».

\subsection{El futuro de la IA en el dispositivo}

Chen Tianqi (陈天奇) compara el panorama actual de la IA con la era de las macrocomputadoras: «Todos usan mainframe y piensan que todo se puede centralize. Pero luego llegó la personal computer. Today la gente sigue acostumbrada a usar el software que tiene en su propia computer; aunque la vision de algo como Chromebook no está mal, la mayoría sigue prefiriendo tener su personal computer.» Los factores impulsores centrales de la implementación en el dispositivo incluyen:

\begin{enumerate}
    \item \textbf{Privacidad (Privacy)}: aunque muchas personas dicen que «no care» (el presentador bromeó diciendo: «si entras a mi ChatGPT y le preguntas por todos mis social security number, seguramente ya los sabe todos»), la necesidad de privacidad existe de manera objetiva
    \item \textbf{Tiempo real}: el presentador, desde la perspectiva del campo de la robótica, señaló que los robots necesitan tomar decisiones en el orden de los milisegundos------«si se va a caminar o a tomar una taza, no hay tiempo de quedarse ahí parado reason 10 segundos». NVIDIA está trabajando recientemente en integrar el chip Thor e implementarlo en robots, lo cual es precisamente un reflejo de esta necesidad
    \item \textbf{Costo}: cuando el preentrenamiento empieza a «start a bit» (el crecimiento de la capacidad del modelo se desacelera), la gente pasa de perseguir el «next level of intelligence» a centrarse en application y cost
    \item \textbf{Especialización}: una vez que el local model alcanza un nivel «suficiente» en un ámbito específico, el panorama cambiará------«no se trata de igualar por completo al modelo más grande en el ámbito general, sino que, si en algunos ámbitos especializados se puede llegar a un nivel suficiente, I think the picture can change»
\end{enumerate}

Chen Tianqi mismo también señaló que, si existiera una IA local con una capacidad casi equivalente a la de GPT, «I would opt to try that way». La variable clave es qué tan fuerte puede llegar a ser realmente la capacidad de los modelos en el dispositivo.

\begin{knowledgebox}{La analogía histórica de Mainframe a Personal Computer}
Esta analogía resulta muy reveladora. La historia de la computación pasó de las macrocomputadoras centralizadas a las computadoras personales distribuidas, y es posible que la IA repita esta trayectoria. Pero no se trata de una simple repetición------la nube y el dispositivo se repartirán las funciones según el escenario, tal como hoy coexisten las aplicaciones web y las aplicaciones locales. La variable central que determina el panorama es el límite de capacidad del modelo local, junto con la ponderación que hace el usuario entre costo, privacidad y tiempo real.
\end{knowledgebox}

\subsection{Resumen de la sección}

MLC LLM refleja la filosofía constante de Chen Tianqi: reducir la barrera de ingeniería mediante tecnología de compilación, para que incluso equipos más pequeños puedan implementar IA en distintos tipos de hardware. El futuro de la IA en el dispositivo depende del equilibrio entre la capacidad del modelo y el costo, pero es muy probable que la tendencia de largo plazo sea un panorama distribuido donde «florezcan cien flores».
